\documentclass[../../main.tex]{subfiles}

\begin{document}

\section{解的存在性}\label{section:ch12-解的存在性}

本节致力于证明温和拟线性初值问题解的存在性定理：
\begin{align}
\begin{cases}
u_{tt}-\Delta u+f(Du,u_t,u)=0 & \text{在 } \mathbb{R}^n\times(0,T] \text{ 中},\\
u=g,\;u_t=h & \text{在 } \mathbb{R}^n\times\{t=0\} \text{ 上}.
\end{cases}
\label{align:ch12-初值问题}
\end{align}
类似 \S7.2，称函数 $u\in L^2(0,T;H^1_{\mathrm{loc}}(\mathbb{R}^n))$，$u'\in L^2(0,T;L^2_{\mathrm{loc}}(\mathbb{R}^n))$，$u''\in L^2(0,T;H^{-1}_{\mathrm{loc}}(\mathbb{R}^n))$ 为初值问题\eqref{align:ch12-初值问题} 的\textbf{弱解}，只要 $u(0)=g$，$u'(0)=h$，$f\coloneqq-f(Du,u',u)\in L^2(0,T;L^2_{\mathrm{loc}}(\mathbb{R}^n))$，且
\begin{align*}
\langle u'',v\rangle+B[u,v]=(f,v)
\end{align*}
对每个 $v\in H^1(\mathbb{R}^n)$ 具有紧支集和 a.e. 时刻 $0\leqslant t\leqslant T$。这里 $B[u,v]\coloneqq\int_{\mathbb{R}^n}Du\cdot Dv\,\mathrm{d}x$。我们总假设 $f(0,0,0)=0$。

\subsection{Lipschitz 非线性}\label{subsection:ch12-Lipschitz非线性}

我们从对非线性性的强假设开始，即
\begin{align}
f:\mathbb{R}^n\times\mathbb{R}\times\mathbb{R}\to\mathbb{R} \text{ 是 Lipschitz 连续的}.
\label{align:ch12-Lipschitz条件}
\end{align}

\begin{theorem}[存在性与唯一性。]\label{theorem:ch12-Lipschitz存在唯一性}
\begin{enumerate}[(i)]
\item 假设 $g\in H^1_{\mathrm{loc}}(\mathbb{R}^n)$，$h\in L^2_{\mathrm{loc}}(\mathbb{R}^n)$。则对每个 $T>0$，存在初值问题\eqref{align:ch12-初值问题} 的唯一弱解 $u$。
\item 若进一步 $g\in H^2_{\mathrm{loc}}(\mathbb{R}^n)$，$h\in H^1_{\mathrm{loc}}(\mathbb{R}^n)$，则
\begin{align*}
u\in L^\infty((0,T);H^2_{\mathrm{loc}}(\mathbb{R}^n)),\;u'\in L^\infty((0,T);H^1_{\mathrm{loc}}(\mathbb{R}^n)),\;u''\in L^\infty((0,T);L^2_{\mathrm{loc}}(\mathbb{R}^n)).
\end{align*}
\end{enumerate}
\end{theorem}

\begin{proof}
1. 首先假设 $T>0$ 充分小。给定 $R\gg 1$，先考虑初边值问题
\begin{align}
\begin{cases}
u_{tt}-\Delta u+f(Du,u_t,u)=0 & \text{在 } B(0,R)\times(0,T] \text{ 中},\\
u=0 & \text{在 } \partial B(0,R)\times[0,T] \text{ 上},\\
u=g,\;u_t=h & \text{在 } B(0,R)\times\{t=0\} \text{ 上}.
\end{cases}
\label{align:ch12-初边值}
\end{align}
暂时也假设 $g\in H^1_0(B(0,R))$。引入空间
\begin{align*}
X\coloneqq\{u\in L^\infty(0,T;H^1_0(B(0,R)))\mid u'\in L^\infty(0,T;L^2(B(0,R)))\},
\end{align*}
范数为
\begin{align*}
\|u\|\coloneqq\operatorname*{ess\,sup}_{0\leqslant t\leqslant T}(\|u(t)\|_{H^1_0(B(0,R))}+\|u'(t)\|_{L^2(B(0,R))}).
\end{align*}
给定 $v\in X$，写 $u=A[v]$ 表示 $u\in X$ 是线性问题
\begin{align}
\begin{cases}
u_{tt}-\Delta u=-f(Dv,v_t,v) & \text{在 } B(0,R)\times(0,T) \text{ 中},\\
u=0 & \text{在 } \partial B(0,R)\times(0,T) \text{ 上},\\
u=g,\;u_t=h & \text{在 } B(0,R)\times\{t=0\} \text{ 上}
\end{cases}
\label{align:ch12-线性问题}
\end{align}
的唯一弱解。此弱解由 \S7.2\cref{theorem:ch07-2-弱解的存在性}--\cref{theorem:ch07-2-改进正则性} 存在。

给定 $\widehat{v}\in X$，写 $\widehat{u}=A[\widehat{v}]$。设 $w\coloneqq u-\widehat{u}$。则 $w$ 是
\begin{align*}
\begin{cases}
w_{tt}-\Delta w=f(D\widehat{v},\widehat{v}_t,\widehat{v})-f(Dv,v_t,v) & \text{在 } B(0,R)\times(0,T) \text{ 中},\\
w=0 & \text{在 } \partial B(0,R)\times(0,T) \text{ 上},\\
w=0,\;w_t=0 & \text{在 } B(0,R)\times\{t=0\} \text{ 上}
\end{cases}
\end{align*}
的唯一弱解。因此 \S7.2 估计给出界
\begin{align*}
\|w\|\leqslant C\|f(D\widehat{v},\widehat{v}',\widehat{v})-f(Dv,v',v)\|_{L^2(0,T;L^2(B(0,R)))}.
\end{align*}
由 $f$ 的 Lipschitz 连续性，
\begin{align*}
\|w\|^2\leqslant C\int_0^T\int_{B(0,R)}|Dv-D\widehat{v}|^2+|v'-\widehat{v}'|^2+|v-\widehat{v}|^2\,\mathrm{d}x\mathrm{d}t\leqslant CT\|v-\widehat{v}\|^2.
\end{align*}
因此
\begin{align*}
\|A[v]-A[\widehat{v}]\|\leqslant\frac{1}{2}\|v-\widehat{v}\|
\end{align*}
只要 $T>0$ 足够小。Banach 不动点定理（\S9.2.1）蕴含存在唯一不动点 $u\in X$，即\eqref{align:ch12-初边值} 的唯一弱解。

2. 特别地可假设 $T<1$。令 $S=R-1$。有限传播速度（\S12.1\cref{theorem:ch12-依赖域}）蕴含柱体 $B(0,S)\times[0,T]$ 内的解仅依赖于限制在 $B(0,S+T)\subset\subset B(0,R)$ 上的初值 $g,h$。因此 $g\in H^1_0(B(0,R))$ 的暂时假设不重要。

对另一个大半径 $\widehat{R}>R>1$ 重复上述构造，由唯一性和有限传播速度，$\widehat{u}\equiv u$ 在 $B(0,R-T)\times[0,T]$ 上。因此可对序列半径 $R_k\to\infty$ 构造解 $u_k$，这些解在 $\mathbb{R}^n\times[0,T]$ 的任意紧子集上存在且一致。这些解对大 $k$ 的公共值确定了\eqref{align:ch12-初值问题} 的唯一弱解。然后将解延拓到时间区间 $[T,2T]$、$[2T,3T]$ 等，构造对所有时间存在的唯一弱解。

3. 选取 $k\in\{1,\cdots,n\}$，令 $\widetilde{u}\coloneqq D_k^hu$ 表示相应的差商（\S5.8.2）。则 $\widetilde{u}$ 是
\begin{align*}
\begin{cases}
\widetilde{u}_{tt}-\Delta\widetilde{u}+\mathbf{b}\cdot D\widetilde{u}+c\widetilde{u}+d u_t=0 & \text{在 } \mathbb{R}^n\times(0,T) \text{ 中},\\
\widetilde{u}=\widetilde{g},\;\widetilde{u}_t=\widetilde{h} & \text{在 } \mathbb{R}^n\times\{t=0\} \text{ 上},
\end{cases}
\end{align*}
的弱解，其中系数由积分中值定理给出。如上估计，对大 $R>0$，
\begin{align*}
\|\widetilde{u}\|=\operatorname*{ess\,sup}_{0\leqslant t\leqslant T}(\|\widetilde{u}(t)\|_{H^1(B(0,R))}+\|\widetilde{u}'(t)\|_{L^2(B(0,R))})\leqslant C.
\end{align*}
上述估计表明 $u\in L^\infty((0,T);H^2_{\mathrm{loc}}(\mathbb{R}^n))$，$u'\in L^\infty((0,T);H^1_{\mathrm{loc}}(\mathbb{R}^n))$。最后利用 PDE 估计 $\int_{B(0,R)}u_{tt}^2\,\mathrm{d}x$，得出 $u''\in L^\infty((0,T);L^2_{\mathrm{loc}}(\mathbb{R}^n))$。
\end{proof}

\subsection{短时间存在性}\label{subsection:ch12-短时间存在性}

再次考虑问题
\begin{align}
\begin{cases}
u_{tt}-\Delta u+f(Du,u_t,u)=0 & \text{在 } \mathbb{R}^n\times(0,T] \text{ 中},\\
u=g,\;u_t=h & \text{在 } \mathbb{R}^n\times\{t=0\} \text{ 上}.
\end{cases}
\label{align:ch12-短时间问题}
\end{align}
现在放弃 $f$ Lipschitz 连续的限制性假设，而仅假设 $f$ 是给定的光滑函数。目标是证明存在至少某个短时间区间 $[0,T]$ 上的唯一解。我们需要初值上更多光滑性，要求 $g\in H^k(\mathbb{R}^n)$，$h\in H^{k-1}(\mathbb{R}^n)$ 对可能大的整数 $k$（依赖于 $n$）。

\begin{theorem}[$H^k$ 的 Sobolev 不等式。]\label{theorem:ch05-Sobolev嵌入}假设函数 $u_1,\cdots,u_m$ 属于 $H^k(\mathbb{R}^n)$，其中 $k>n/2$。
\begin{enumerate}[(i)]
\item 若 $|\beta_1|+\cdots+|\beta_m|\leqslant k$，则
\begin{align}
\|D^{\beta_1}u_1\cdots D^{\beta_m}u_m\|_{L^2(\mathbb{R}^n)}\leqslant C\prod_{j=1}^m\|u_j\|_{H^k(\mathbb{R}^n)}
\label{align:ch12-Hk乘积估计}
\end{align}
对常数 $C=C(n,m,k)$。
\item 设 $f:\mathbb{R}^m\to\mathbb{R}$ 是光滑函数满足 $f(0)=0$。则 $f(u_1,\cdots,u_m)\in H^k(\mathbb{R}^n)$ 且
\begin{align}
\|f(u_1,\cdots,u_m)\|_{H^k(\mathbb{R}^n)}\leqslant\Phi(\|u_1\|_{H^k(\mathbb{R}^n)},\cdots,\|u_m\|_{H^k(\mathbb{R}^n)}),
\label{align:ch12-Hk复合估计}
\end{align}
其中 $\Phi$ 是连续函数，在每个变量中非降，仅依赖于 $f,k,n,m$。
\end{enumerate}
\end{theorem}

\begin{proof}
1. \eqref{align:ch12-Hk乘积估计} 的证明留作读者（见\cref{exercise:ch12-SobolevHk}）。

2. 设 $0<|\alpha|\leqslant k$。则 $D^\alpha f(u_1,\cdots,u_m)$ 可写为有限和的项 $AD^{\beta_1}u_{j_1}\cdots D^{\beta_l}u_{j_l}$，其中 $A$ 依赖于 $f$ 的至多 $k$ 阶偏导数在 $(u_1,\cdots,u_m)$ 处的值，$l\leqslant k$，$0\leqslant|\beta_j|\leqslant|\alpha|$，$\beta_1+\cdots+\beta_l=\alpha$。回忆 \S5.6.3\cref{theorem:ch05-一般Sobolev不等式}，$k>n/2$ 蕴含估计 $\|u\|_{L^\infty}\leqslant C\|u\|_{H^k}$。因此 $\|A\|_{L^\infty}$ 被仅依赖于 $f$ 和 $\|u_j\|_{H^k}$（$j=1,\cdots,m$）的项界住。由\eqref{align:ch12-Hk乘积估计}，
\begin{align*}
\|AD^{\beta_1}u_{j_1}\cdots D^{\beta_l}u_{j_l}\|_{L^2}
\end{align*}
被仅涉及 $\|u_j\|_{H^k}$（$j=1,\cdots,m$）的表达式界住。由于 $f(0)=0$，可类似估计 $\|f(u_1,\cdots,u_m)\|_{L^2}$。这建立了\eqref{align:ch12-Hk复合估计}。
\end{proof}

\begin{theorem}[短时间存在性。]\label{theorem:ch12-短时间存在性}假设 $f:\mathbb{R}^n\times\mathbb{R}\times\mathbb{R}\to\mathbb{R}$ 光滑，$f(0,0,0)=0$。又设 $g\in H^k(\mathbb{R}^n)$，$h\in H^{k-1}(\mathbb{R}^n)$，其中 $k>n/2+1$。则存在时刻 $T>0$ 使得初值问题
\begin{align}
\begin{cases}
u_{tt}-\Delta u+f(Du,u_t,u)=0 & \text{在 } \mathbb{R}^n\times(0,T] \text{ 中},\\
u=g,\;u_t=h & \text{在 } \mathbb{R}^n\times\{t=0\} \text{ 上}
\end{cases}
\label{align:ch12-短时间初值}
\end{align}
具有唯一弱解 $u$，且
\begin{align*}
u\in L^\infty(0,T;H^k(\mathbb{R}^n)),\;u'\in L^\infty(0,T;H^{k-1}(\mathbb{R}^n)).
\end{align*}
证明提供的存在时间 $T$ 以复杂方式依赖于 $f$ 和 $\|g\|_{H^k}$、$\|h\|_{H^{k-1}}$，若 $\|g\|_{H^k}$ 和 $\|h\|_{H^{k-1}}$ 大则 $T$ 可非常短。
\end{theorem}

\begin{proof}
1. 设
\begin{align*}
X\coloneqq\{u\in L^\infty(0,T;H^1(\mathbb{R}^n))\mid u'\in L^\infty(0,T;L^2(\mathbb{R}^n))\}
\end{align*}
范数为
\begin{align*}
\|u\|\coloneqq\operatorname*{ess\,sup}_{0\leqslant t\leqslant T}(\|u(t)\|_{H^1(\mathbb{R}^n)}+\|u'(t)\|_{L^2(\mathbb{R}^n)}).
\end{align*}
引入更强范数
\begin{align*}
\lll u\rrr\coloneqq\operatorname*{ess\,sup}_{0\leqslant t\leqslant T}(\|u(t)\|_{H^k(\mathbb{R}^n)}+\|u'(t)\|_{H^{k-1}(\mathbb{R}^n)}).
\end{align*}
对 $\lambda>0$ 定义
\begin{align*}
X_\lambda\coloneqq\{u\in X\mid\lll u\rrr\leqslant\lambda,\;u(0)=g,\;u'(0)=h\}.
\end{align*}
若 $v\in X_\lambda$，写 $u=A[v]$ 表示 $u$ 求解线性初值问题
\begin{align}
\begin{cases}
u_{tt}-\Delta u=-f(Dv,v_t,v) & \text{在 } \mathbb{R}^n\times(0,T] \text{ 中},\\
u=g,\;u_t=h & \text{在 } \mathbb{R}^n\times\{t=0\} \text{ 上}.
\end{cases}
\label{align:ch12-线性短时间}
\end{align}
定义 $E_k(t)\coloneqq\|u(t)\|_{H^k}^2+\|u'(t)\|_{H^{k-1}}^2$ 和 $F_k(t)\coloneqq\|v(t)\|_{H^k}^2+\|v'(t)\|_{H^{k-1}}^2$。

2. 我们断言有估计
\begin{align}
E_k(t)\leqslant E_k(0)+C\int_0^t\Psi(F_k(s))\,\mathrm{d}s \quad (0\leqslant t\leqslant T)
\label{align:ch12-Ek估计}
\end{align}
对某连续单调函数 $\Psi$ 仅依赖于 $n,k,f$。

为证明\eqref{align:ch12-Ek估计}，取 $|\alpha|\leqslant k-1$，对 PDE \eqref{align:ch12-线性短时间} 作用 $D^\alpha$，令 $w\coloneqq D^\alpha u$，得 $w_{tt}-\Delta w=-D^\alpha(f(Dv,v_t,v))$。利用\eqref{align:ch12-Hk复合估计}（以 $k-1$ 代替 $k$，$m=n+2$），计算
\begin{align*}
\frac{\mathrm{d}}{\mathrm{d}t}\int_{\mathbb{R}^n}(w_t^2+|Dw|^2)\,\mathrm{d}x
&=2\int_{\mathbb{R}^n}w_t(w_{tt}-\Delta w)\,\mathrm{d}x
=-2\int_{\mathbb{R}^n}w_tD^\alpha f(Dv,v_t,v)\,\mathrm{d}x\\
&\leqslant\int_{\mathbb{R}^n}w_t^2+|D^\alpha f(Dv,v_t,v)|^2\,\mathrm{d}x
\leqslant\int_{\mathbb{R}^n}w_t^2\,\mathrm{d}x+\Psi(F_k(t)).
\end{align*}
应用 Gronwall 不等式并对所有 $|\alpha|\leqslant k-1$ 求和，推出\eqref{align:ch12-Ek估计}。

3. 我们断言若 $\lambda>0$ 足够大且 $T>0$ 足够小，则 $A:X_\lambda\to X_\lambda$。由\eqref{align:ch12-Ek估计}，
\begin{align*}
\lll u\rrr^2\leqslant\|g\|_{H^k}^2+\|h\|_{H^{k-1}}^2+CT\Psi(\lll v\rrr^2).
\end{align*}
设 $\lambda^2\coloneqq 2(\|g\|_{H^k}^2+\|h\|_{H^{k-1}}^2)$，然后固定 $T>0$ 足够小使得 $CT\Psi(\lambda^2)\leqslant\lambda^2/2$。则 $\lll A[v]\rrr^2\leqslant\lambda^2$。

4. 接下来断言若 $\lambda$ 足够大且 $T$ 足够小，则
\begin{align}
\|A[v]-A[\widehat{v}]\|\leqslant\frac{1}{2}\|v-\widehat{v}\|
\label{align:ch12-压缩}
\end{align}
对所有 $v,\widehat{v}\in X_\lambda$。写 $u=A[v]$，$\widehat{u}=A[\widehat{v}]$，设 $w\coloneqq u-\widehat{u}$。则
\begin{align*}
\frac{\mathrm{d}}{\mathrm{d}t}\int_{\mathbb{R}^n}|Dw|^2+w_t^2+w^2\,\mathrm{d}x
&=2\int_{\mathbb{R}^n}w_t(f(D\widehat{v},\widehat{v}_t,\widehat{v})-f(Dv,v_t,v)+w)\,\mathrm{d}x\\
&\leqslant\int_{\mathbb{R}^n}w_t^2+w^2\,\mathrm{d}x
+C\int_{\mathbb{R}^n}|D\widehat{v}-Dv|^2+|\widehat{v}_t-v_t|^2+|\widehat{v}-v|^2\,\mathrm{d}x,
\end{align*}
常数 $C$ 依赖于 $\|D\widehat{v},\widehat{v}_t,\widehat{v},Dv,v_t,v\|_{L^\infty}$。此量有界，因 $v,\widehat{v}\in X_\lambda$ 且 $k>n/2+1$。应用 Gronwall 不等式，
\begin{align*}
\max_{0\leqslant t\leqslant T}\int|Dw|^2+|w_t|^2+w^2\,\mathrm{d}x\leqslant C\int_0^T\int_{\mathbb{R}^n}|Dv-D\widehat{v}|^2+|v_t-\widehat{v}_t|^2+|v-\widehat{v}|^2\,\mathrm{d}x\mathrm{d}t\leqslant CT\|v-\widehat{v}\|^2.
\end{align*}
选取 $T$ 足够小以确保\eqref{align:ch12-压缩}。

5. 选取任意 $u_0\in X_\lambda$。由 \S9.2.1 Banach 定理的证明，归纳定义 $u_{k+1}\coloneqq A[u_k]$，则 $u_k\to u$ 在 $X$ 中且 $A[u]=u$。此外因 $\lll u_k\rrr\leqslant\lambda$，$u\in X_\lambda$。唯一性由\eqref{align:ch12-压缩} 得。
\end{proof}

注意此证明的策略：我们证明对小 $T>0$，算子 $A$ 在较弱范数 $\|\cdot\|$ 下是严格压缩的，同时在较强范数 $\lll\cdot\rrr\rrr$ 下保持某些估计。因此 Banach 定理的迭代格式提供了在 $X$ 中收敛到不动点的序列，该不动点实际上位于更好的空间 $X_\lambda$ 中。我们不必证明 $A$ 在较强范数下是严格压缩的。

\end{document}